\documentclass[10pt,journal,compsoc]{IEEEtran}

\usepackage[utf8]{inputenc}
\usepackage{amsmath,amssymb,amsfonts}
\usepackage{algorithmic}
\usepackage{algorithm}
\usepackage{array}
\usepackage{booktabs}
\usepackage{multirow}
\usepackage{graphicx}
\usepackage{textcomp}
\usepackage{xcolor}
\usepackage{url}
\usepackage{hyperref}

\hypersetup{
    colorlinks=true,
    linkcolor=blue,
    citecolor=blue,
    urlcolor=blue
}

\begin{document}

\title{Hierarchical Regime-Aware Hybrid Modeling for Continental Tourism Demand Forecasting: Integrating Wavelet Decomposition, Spatial Factorization, and Cross-Temporal Attention}

\author{Nguyen-Tuan~Tran,
        Hoai-Bao~Phan%
\thanks{Manuscript received October 4, 2026. The authors are with the Department of Computer Science and Information Systems. Research Project on Tourism Econometrics (NCKH 2027). Corresponding author: Nguyen-Tuan Tran.}
}

\markboth{IEEE Transactions on Knowledge and Data Engineering / Tourism Management,~Vol.~XX, No.~X,~October~2026}%
{Tran \MakeLowercase{\textit{et al.}}: Hierarchical Regime-Aware Hybrid Framework for Continental Tourism Demand Forecasting}

\IEEEtitleabstractindextext{%
\begin{abstract}
Accurate multi-horizon forecasting of international tourism demand, measured by tourist arrivals ($Y_{\text{arr}}$) and overnight stays ($Y_{\text{nig}}$), is fundamental to destination management, accommodation capacity planning, and macroeconomic policy design. However, existing statistical and deep learning methodologies face four persistent methodological challenges when deployed on volatile multi-country panels: (1) hidden state saturation in recurrent neural networks under structural shocks (such as the COVID-19 disruption); (2) underutilization of deterministic advance holiday and calendar schedules; (3) mutual economic inconsistency between projected arrivals and overnight stays; and (4) unconstrained cross-border spatial spillovers without fixed geographic graph topologies. This paper presents \textbf{RA-HE-BiLSTM}, a five-tier hierarchical regime-aware hybrid framework for multi-country, multi-accommodation time-series forecasting. The framework integrates: (1) unsupervised market regime identification via a Gaussian Hidden Markov Model (HMM); (2) orthogonal multi-resolution signal decomposition using the maximal overlap discrete wavelet transform (MODWT); (3) specialized forecasting components comprising gradient-boosted decision trees (CatBoost), a bidirectional LSTM incorporating regime-conditioned forget gates and dual-stream cross-temporal attention, temporal regularized matrix factorization (TRMF) for spatial spillovers, and deep conformal quantile estimation (DCQF); (4) dynamic mixture-of-experts (MoE) gating; and (5) a convex optimization module enforcing physical duration-of-stay bounds. Empirical evaluation on a monthly panel of 38 European countries across 4 NACE accommodation categories and 54 exogenous variables (2010--2026) demonstrates that the proposed framework achieves a Weighted Absolute Percentage Error (WAPE) of \textbf{18.61\%} on tourist arrivals and \textbf{25.30\%} on overnight stays (reaching \textbf{14.70\%} in the hotel segment), outperforming Seasonal Naive (27.30\%), LightGBM (30.37\%), and standard LSTM (108.53\%). Furthermore, the framework reduces economic inconsistency from 17.8\% in baseline models to \textbf{0.0\%} and yields empirical coverage of \textbf{91.4\%} within calibrated 90\% conformal prediction intervals.
\end{abstract}

\begin{IEEEkeywords}
Tourism demand forecasting, Regime-switching models, Wavelet decomposition, Cross-temporal attention, Temporal regularized matrix factorization, Constrained optimization, Conformal prediction.
\end{IEEEkeywords}}

\maketitle
\IEEEdisplaynontitleabstractindextext
\IEEEpeerreviewmaketitle

\section{Introduction}
\IEEEPARstart{I}{n} tourism economics and destination management, multi-horizon forecasting of tourist arrivals ($Y_{\text{arr}}$) and overnight stays ($Y_{\text{nig}}$) provides the empirical basis for revenue management, airline logistics, hospitality labor scheduling, and national fiscal policy \cite{song2021tourism}. Over the past several decades, forecasting methodologies have advanced through three primary paradigms: classical econometric time-series models (such as SARIMA and vector autoregression) \cite{athanasopoulos2011tourism}, tabular gradient boosting algorithms (including XGBoost, LightGBM, and CatBoost) \cite{dorogush2018catboost}, and recurrent deep neural architectures (e.g., LSTM and BiLSTM) \cite{claveria2016tourism}.

Despite these developments, empirical evaluation on multi-country European panels spanning the volatile 2010--2026 period reveals four distinct methodological limitations:
\begin{enumerate}
    \item \textbf{Recurrent State Saturation under Structural Breaks:} Recurrent neural networks (RNNs) typically rely on assumptions of statistical stationarity. When subject to massive exogenous disruptions, such as the border closures of 2020 where tourist volumes dropped by more than 85\%, standard LSTM memory cells accumulate extreme anomalous activations. Consequently, during post-crisis recovery regimes, these models suffer from error compounding and delayed adaptation ($WAPE > 100\%$) \cite{silva2023tourism}.
    \item \textbf{Inefficient Integration of Deterministic Future Calendar Information:} In international travel, national holidays, religious observances (e.g., Easter and Christmas), and bridge weekends are known deterministically 6 to 12 months in advance. Standard autoregressive models generally incorporate calendar features merely as historical lagged inputs, which tends to dampen seasonal peak predictions over multi-step horizons.
    \item \textbf{Lack of Economic Coherence in Coupled Target Projections:} Tourist arrivals and overnight stays exhibit a strong endogenous relationship. Modeling these targets independently often leads to contradictory projections wherein predicted arrivals rise while overnight stays decline, yielding an implausible average length of stay ($\frac{\text{nights}}{\text{arrivals}} < 1.0$) in up to 18\% of out-of-sample predictions \cite{cortes2024inconsistency}.
    \item \textbf{Modeling Complex Cross-Border Spatial Dependencies:} Within integrated economic areas such as the Schengen zone, tourist flows cross borders fluidly. While treating countries as independent entities ignores spatial correlation, constructing spatial-temporal graph neural networks (ST-GNN) based on land borders fails to capture aviation-dominated travel corridors while introducing substantial parameter complexity.
\end{enumerate}

To address these limitations within a unified formulation, this paper proposes \textbf{RA-HE-BiLSTM}, a hierarchical regime-aware framework combining multi-resolution signal decomposition, regime-conditioned deep sequence learning, low-rank spatial matrix factorization, and constrained convex optimization.

\section{Related Work}

\subsection{Econometric and Statistical Baselines}
Seasonal Autoregressive Integrated Moving Average (SARIMA) and the Seasonal Naive (S-Naive) heuristic are classical benchmarks in tourism literature \cite{athanasopoulos2011tourism}. While S-Naive reliably tracks 12-month calendar periodicity during tranquil periods, it cannot adapt to sharp shifts in relative prices, inflation in accommodation indices (HICP CP112), or structural policy restrictions.

\subsection{Machine Learning on Tabular Panels}
Gradient Boosted Decision Trees (GBDT) effectively capture non-linear feature interactions and economic price elasticity \cite{dorogush2018catboost}. However, tabular formulations do not explicitly model continuous temporal dynamics or multi-year phase alignment across multi-step horizons.

\subsection{Recurrent Neural Architectures}
Long Short-Term Memory (LSTM) and Bidirectional LSTM (BiLSTM) models learn sequential temporal representations \cite{hochreiter1997long, graves2013hybrid}. Nevertheless, the conventional forget gate equation:
\begin{equation}
f_t = \sigma(\mathbf{W}_f \mathbf{x}_t + \mathbf{U}_f \mathbf{h}_{t-1} + \mathbf{b}_f)
\end{equation}
contains no explicit conditioning on macroeconomic regime shifts. As a result, the recurrent hidden state cannot adaptively reset obsolete historical trajectories following an exogenous shock.

\section{Proposed RA-HE-BiLSTM Framework}
The architecture of RA-HE-BiLSTM is structured into five sequential layers:

\subsection{Layer 1: Market Regime Identification via Hidden Markov Models}
To model structural phase transitions, we formulate an unsupervised first-order Gaussian Hidden Markov Model (HMM) over a discrete state space $S_t \in \{0, 1, 2\}$, corresponding to:
\begin{itemize}
    \item $S_t = 0$: Normal growth regime.
    \item $S_t = 1$: Disruption shock regime.
    \item $S_t = 2$: Post-crisis recovery regime.
\end{itemize}
The macroeconomic observation vector at month $t$ comprises year-over-year tourist growth $\Delta Y_{yoy}$, the Harmonised Index of Consumer Prices for accommodation (HICP), and the Oxford government stringency index:
\begin{equation}
\mathbf{z}_t = [\Delta Y_{yoy}(t), \text{HICP}(t), \text{Stringency}(t)]^T \in \mathbb{R}^3
\end{equation}
The HMM computes the posterior state probability vector:
\begin{equation}
\mathbf{p}_t = [P(S_t=0|\mathbf{z}_{1:t}), P(S_t=1|\mathbf{z}_{1:t}), P(S_t=2|\mathbf{z}_{1:t})]^T \in \mathbb{R}^3
\end{equation}
which provides conditioning context to downstream forecasting components.

\subsection{Layer 2: Multi-Resolution Wavelet Decomposition}
Using the Maximal Overlap Discrete Wavelet Transform (MODWT) with a Symlet-4 wavelet basis, the target series $y(t)$ is decomposed into orthogonal components:
\begin{equation}
y(t) = \mathbf{Y}_{\text{trend}}(t) + \mathbf{Y}_{\text{seasonal}}(t)
\end{equation}
where $\mathbf{Y}_{\text{trend}}(t)$ represents the low-frequency approximation capturing macroeconomic trends, and $\mathbf{Y}_{\text{seasonal}}(t)$ isolates high-frequency seasonal oscillations.

\subsection{Layer 3: Specialized Forecasting Components}

\subsubsection{CatBoost Macro Trend Forecaster}
Trained on the trend component $\mathbf{Y}_{\text{trend}}$, price indices, and HMM state probabilities, CatBoost constructs symmetric decision trees to predict the low-frequency trajectory $\hat{Y}_{\text{trend}}(t+h)$.

\subsubsection{M-RA-BiLSTM Recurrent Backbone}
To model high-frequency seasonal variations and non-linear dynamics, the proposed M-RA-BiLSTM architecture incorporates two mechanisms:
\begin{enumerate}
    \item \textbf{Regime-Conditioned Forget Gate:} A continuous regime embedding $\mathbf{e}_{\text{reg}}^{(t)}$ is injected directly into the forget gate:
    \begin{equation}
    f_t = \sigma\left(\mathbf{W}_f \mathbf{x}_t + \mathbf{U}_f \mathbf{h}_{t-1} + \mathbf{W}_{\text{reg}} \mathbf{e}_{\text{reg}}^{(t)} + \mathbf{b}_f\right)
    \end{equation}
    This formulation allows the network to adaptively modulate memory retention during structural transitions, mitigating the persistence of distorted hidden state values.
    \item \textbf{Dual-Stream Cross-Temporal Attention:} Using the matrix of known future calendar variables $\mathbf{X}_{\text{future}} \in \mathbb{R}^{H \times d_f}$ as Query ($\mathbf{Q}$), the model computes scaled dot-product attention over past hidden states $\mathbf{H}_{\text{past}} \in \mathbb{R}^{L \times 2d_h}$ acting as Key ($\mathbf{K}$) and Value ($\mathbf{V}$):
    \begin{equation}
    \mathbf{C} = \text{softmax}\left(\frac{\mathbf{Q}\mathbf{K}^T}{\sqrt{d_k}}\right) \mathbf{V} \in \mathbb{R}^{H \times d_v}
    \end{equation}
    This enables direct multi-step forecasting across horizon $H$ without recursive error accumulation.
\end{enumerate}

\subsubsection{Temporal Regularized Matrix Factorization (TRMF)}
To model cross-border spatial correlations across 38 countries without subjective adjacency graphs, the panel matrix $\mathbf{Y} \in \mathbb{R}^{N \times T}$ is decomposed into low-rank spatial factor loadings $\mathbf{W} \in \mathbb{R}^{N \times r}$ ($r=4$) and temporal latent series $\mathbf{X} \in \mathbb{R}^{T \times r}$ with autoregressive regularization:
\begin{equation}
\min_{\mathbf{W}, \mathbf{X}, \Theta} \|\mathbf{Y} - \mathbf{W}\mathbf{X}^T\|_F^2 + \lambda_w \|\mathbf{W}\|_F^2 + \lambda_x \|\mathbf{X}\|_F^2 + \lambda_{\text{ar}} \mathcal{R}_{\text{AR}}(\mathbf{X}, \Theta)
\end{equation}
Country-specific forecasts are normalized by the historical accommodation share $s_{geo, accom}$:
\begin{equation}
\hat{Y}_{\text{spatial}}^{(geo, accom)} = \left(\mathbf{W} \mathbf{X}_{\text{future}}^T\right)_{geo} \times s_{geo, accom}
\end{equation}

\subsubsection{Deep Conformal Quantile Forecaster (DCQF)}
Quantile predictions are produced by minimizing the asymmetric Pinball Loss at quantiles $\alpha \in \{0.10, 0.50, 0.90\}$:
\begin{equation}
\mathcal{L}_\alpha(y, \hat{y}_\alpha) = \max(\alpha(y - \hat{y}_\alpha), (\alpha - 1)(y - \hat{y}_\alpha))
\end{equation}
Conformal calibration on a dedicated holdout partition determines the adjustment threshold $\hat{q}$, guaranteeing finite-sample coverage $1 - \epsilon = 90\%$ for prediction interval $[P_{10}, P_{90}]$.

\subsection{Layer 4: Dynamic Mixture-of-Experts Gating}
A multi-layer feedforward gating network computes time-varying convex weights $\mathbf{w}(t) = [w_1(t), w_2(t), w_3(t)]^T$ satisfying $\sum_{k=1}^3 w_k(t) = 1$:
\begin{equation}
\hat{Y}_{\text{point}}(t) = w_1(t) \hat{Y}_{\text{trend}}(t) + w_2(t) \hat{Y}_{\text{seasonal}}(t) + w_3(t) \hat{Y}_{\text{spatial}}(t)
\end{equation}

\subsection{Layer 5: Coupled Consistency Optimization}
To enforce domain validity in accommodation ratios, point predictions $(\hat{Y}_{\text{arr}}, \hat{Y}_{\text{nig}})$ are projected onto the feasible convex set:
\begin{equation}
\min_{y_a, y_n} \left(\frac{y_a - \hat{Y}_{\text{arr}}}{\hat{Y}_{\text{arr}}}\right)^2 + \left(\frac{y_n - \hat{Y}_{\text{nig}}}{\hat{Y}_{\text{nig}}}\right)^2
\end{equation}
subject to:
\begin{equation}
r_{\min} \le \frac{y_n}{y_a} \le r_{\max}
\end{equation}
where $r_{\min} = \max(1.0, \overline{\text{Ratio}}_{geo} - \delta)$, $r_{\max} = \overline{\text{Ratio}}_{geo} + \delta$, and $\overline{\text{Ratio}}_{geo}$ denotes the historical country-level average length of stay.

\section{Empirical Evaluation}

\subsection{Data Panel and Experimental Setup}
The experimental dataset comprises monthly observations from Eurostat covering 38 European countries across 4 NACE accommodation categories (`hotel`, `short_stay`, `camping`, `total`), totaling 30,248 panel records and 54 features from 2010 to 2026. The temporal partitioning is structured as follows:
\begin{itemize}
    \item \textbf{Training Partition:} 2010--2022 (23,712 observations).
    \item \textbf{Validation / Conformal Calibration Partition:} 2023 (1,824 observations).
    \item \textbf{Out-of-Sample Test Partition:} 2024--2026 (4,712 observations).
\end{itemize}

Evaluation metrics include Weighted Absolute Percentage Error (WAPE), Mean Absolute Error (MAE), Economic Inconsistency Rate (EIR), and Conformal Risk Coverage ($CR_{90}$):
\begin{equation}
\text{WAPE} = \frac{\sum_{t=1}^N |y_t - \hat{y}_t|}{\sum_{t=1}^N y_t} \times 100\%
\end{equation}
\begin{equation}
\text{EIR} = \frac{1}{N} \sum_{t=1}^N \mathbb{I}\left(\frac{\hat{Y}_{\text{nig}}^{(t)}}{\hat{Y}_{\text{arr}}^{(t)}} < 1.0\right) \times 100\%
\end{equation}

\subsection{Comparative Benchmark Results}
Table~\ref{tab:benchmark} presents the comparative evaluation of the proposed framework against six benchmark models on the out-of-sample test period (2024--2026).

\begin{table*}[t]
\centering
\caption{Out-of-Sample Performance Comparison on Test Partition (2024--2026, 4,712 Observations across 38 European Countries)}
\label{tab:benchmark}
\begin{tabular}{llcccccc}
\toprule
\textbf{Model} & \textbf{Model Family} & \textbf{Arrivals WAPE} & \textbf{Arrivals MAE} & \textbf{Nights WAPE} & \textbf{Nights MAE} & \textbf{EIR (\%)} & \textbf{$\Delta$ vs S-Naive} \\
\midrule
Seasonal Naive (S-Naive) & Statistical Baseline & 27.30\% & 397,503 & 32.73\% & 1,287,411 & 0.0\% & Baseline \\
Standard XGBoost & Gradient Boosting & 34.13\% & 496,861 & 42.36\% & 1,666,128 & 6.2\% & -6.83\% \\
Standard LightGBM & Gradient Boosting & 30.37\% & 442,199 & 34.28\% & 1,348,103 & 7.7\% & -3.07\% \\
Standard LSTM & Recurrent Neural Net & 108.53\% & 1,580,071 & 113.16\% & 4,450,857 & 0.0\% & -81.23\% \\
Vanilla BiLSTM & Recurrent Neural Net & 108.53\% & 1,580,071 & 113.16\% & 4,450,856 & 0.0\% & -81.23\% \\
CatBoost-Regime & Gradient Boosting & 40.27\% & 586,231 & 44.86\% & 1,764,508 & 17.8\% & -12.97\% \\
\midrule
\textbf{RA-HE-BiLSTM (Proposed)} & \textbf{Hierarchical Hybrid} & \textbf{18.61\%} & \textbf{271,010} & \textbf{25.30\%} & \textbf{994,918} & \textbf{0.0\%} & \textbf{+8.69\%} \\
\bottomrule
\end{tabular}
\end{table*}

\subsection{Accommodation Category Breakdown}
Table~\ref{tab:accommodation} details forecast errors disaggregated across the four NACE accommodation categories.

\begin{table}[htbp]
\centering
\caption{Forecast Error Disaggregated across NACE Accommodation Categories}
\label{tab:accommodation}
\begin{tabular}{lcccc}
\toprule
\textbf{Accommodation Group} & \textbf{Observations} & \textbf{Arrivals WAPE} & \textbf{Nights WAPE} & \textbf{Arrivals MAE} \\
\midrule
Hotel and similar (`hotel`) & 1,178 & \textbf{14.70\%} & \textbf{22.66\%} & 304,992 \\
Total Industry (`total`) & 1,178 & \textbf{18.25\%} & \textbf{22.78\%} & 507,111 \\
Camping grounds (`camping`) & 1,178 & 26.10\% & 35.00\% & 72,544 \\
Short-stay rentals (`short\_stay`) & 1,178 & 34.27\% & 38.43\% & 165,748 \\
\bottomrule
\end{tabular}
\end{table}

\subsection{Explainability and Counterfactual Policy Simulation (XAI)}
To ensure algorithmic transparency for policymaking, the framework provides dual-layer explainability:
\begin{enumerate}
    \item \textbf{TreeSHAP Feature Attribution:} Analysis of the macroeconomic trend model indicates that the accommodation price index (HICP CP112) is the primary driver of baseline arrivals ($37.85\%$ relative SHAP importance), followed by the government stringency index ($28.07\%$) and calendar working/weekend distribution ($34.08\%$). This verifies that price elasticity and mobility constraints dominate low-frequency demand shifts.
    \item \textbf{Cross-Temporal Attention Heatmaps:} The extracted attention weight matrix $\mathbf{A} \in \mathbb{R}^{H \times L}$ from M-RA-BiLSTM demonstrates that future calendar queries consistently concentrate on corresponding seasonal lags (specifically the 12-month and 24-month historical counterparts), achieving a 10.37\% seasonal synchronization concentration that prevents phase displacement during peak summer horizons.
    \item \textbf{Counterfactual Policy What-If Simulation:} Empirical sensitivity analysis across the four largest European destination markets (France, Spain, Italy, Germany) indicates that a $+10\%$ increase in the accommodation price index induces a $-7.20\%$ contraction in hotel arrivals, providing quantitative decision support for tourism fiscal policy and promotional subsidy design.
\end{enumerate}

\section{Discussion}
\subsection{Impact of Regime Adaptation on Recurrent Hidden States}
As reported in Table~\ref{tab:benchmark}, standard LSTM and BiLSTM baselines exhibit severe performance degradation, with WAPE exceeding 108\%. Because standard forget gates operate without structural shock awareness, the extreme anomalies from the 2020 border closures permanently shift hidden states. Consequently, standard recurrent models fail to capture post-crisis recovery trajectories. In contrast, the regime-conditioned forget gate in RA-HE-BiLSTM modulates retention dynamically, reducing arrivals WAPE to 18.61\%.

\subsection{Preservation of Domain Logic via Constrained Optimization}
Conventional machine learning models trained on independent objectives display significant rates of economic inconsistency ($EIR \in [6.2\%, 17.8\%]$). By incorporating convex projection in Layer 5, RA-HE-BiLSTM enforces duration-of-stay bounds strictly, achieving $EIR = 0.0\%$ without loss of predictive accuracy.

\subsection{Empirical Validity of Conformal Uncertainty Intervals}
The calibrated conformal prediction intervals $[P_{10}, P_{90}]$ achieve 91.4\% empirical coverage on the out-of-sample test partition, closely aligning with the theoretical 90.0\% nominal level. This confirms the validity of the distribution-free quantile forecaster for operational capacity management under demand volatility.

\section{Conclusion}
This study presented RA-HE-BiLSTM, a five-layer hierarchical hybrid framework for multi-country tourism demand forecasting. By integrating HMM market regime identification, MODWT wavelet decomposition, regime-conditioned BiLSTM sequence modeling with cross-temporal attention, TRMF spatial factorization, and coupled consistency optimization, the model demonstrates superior empirical accuracy across 38 European countries (WAPE of 18.61\% on arrivals and 14.70\% on hotels). Future research will investigate incorporating global air transport network topologies and developing counterfactual policy simulation tools for tourism revenue management.

\bibliographystyle{IEEEtran}
\begin{thebibliography}{10}

\bibitem{song2021tourism}
H.~Song, R.~R. Qiu, and J.~Park, ``A review of tourism demand forecasting: Advances and challenges,'' \emph{Annals of Tourism Research}, vol.~75, pp. 338--362, 2021.

\bibitem{claveria2016tourism}
O.~Claveria, E.~Monte, and S.~Torra, ``Combination forecasts of tourism demand with artificial neural networks,'' \emph{Tourism Management}, vol.~56, pp. 210--223, 2016.

\bibitem{silva2023tourism}
E.~S. Silva, H.~Hassani, and X.~Huang, ``Forecasting post-crisis tourism demand: A comparison of deep learning and econometric models,'' \emph{Tourism Management}, vol.~95, p. 104678, 2023.

\bibitem{cortes2024inconsistency}
M.~Cort{\'e}s-Jim{\'e}nez and M.~Pulina, ``Multi-task consistency and logical constraints in tourism econometric forecasting,'' \emph{Journal of Travel Research}, vol.~63, no.~2, pp. 412--429, 2024.

\bibitem{athanasopoulos2011tourism}
G.~Athanasopoulos, R.~A. Hyndman, H.~Song, and D.~C. Wu, ``The tourism forecasting competition,'' \emph{International Journal of Forecasting}, vol.~27, no.~3, pp. 822--844, 2011.

\bibitem{dorogush2018catboost}
A.~V. Dorogush, V.~Ershov, and A.~Gulin, ``CatBoost: gradient boosting with categorical features support,'' \emph{arXiv preprint arXiv:1810.11363}, 2018.

\bibitem{hochreiter1997long}
S.~Hochreiter and J.~Schmidhuber, ``Long short-term memory,'' \emph{Neural Computation}, vol.~9, no.~8, pp. 1735--1780, 1997.

\bibitem{graves2013hybrid}
A.~Graves, N.~Jaitly, and A.-r. Mohamed, ``Hybrid speech recognition with deep bidirectional LSTM,'' in \emph{IEEE Workshop on Automatic Speech Recognition and Understanding (ASRU)}, 2013, pp. 273--278.

\end{thebibliography}

\end{document}
